% Общий стиль PDF курса «AI-инженер» (подключается ко всем PDF через _quarto.yml)

\usepackage{xcolor}
\definecolor{mlaccent}{HTML}{3949AB}
\definecolor{mlmuted}{HTML}{6B7280}

% Заголовки — в фирменном цвете
\addtokomafont{disposition}{\color{mlaccent}}
\setkomafont{subtitle}{\normalfont\sffamily\large\color{mlmuted}}
\setkomafont{subject}{\normalfont\sffamily\small\color{mlmuted}}

% Строка над названием документа
\subject{AI-инженер: от устройства LLM до мультиагентных систем}

% Колонтитулы
\usepackage[automark]{scrlayer-scrpage}
\clearpairofpagestyles
\ihead{\small\sffamily\color{mlmuted}AI-инженер}
\ohead{\small\sffamily\color{mlmuted}\headmark}
\cfoot*{\small\sffamily\color{mlmuted}\pagemark}
\pagestyle{scrheadings}

% Типографика: меньше «дыр» в строках, аккуратные таблицы
\setlength{\emergencystretch}{3em}
\frenchspacing
\renewcommand{\arraystretch}{1.25}

% В PT Serif нет стрелок: делаем их активными символами и берём из математического шрифта
% (то же, что пакет newunicodechar, но без зависимостей)
\catcode`→=\active \protected\def→{\ensuremath{\rightarrow}}
\catcode`←=\active \protected\def←{\ensuremath{\leftarrow}}
\catcode`↔=\active \protected\def↔{\ensuremath{\leftrightarrow}}
\catcode`⇒=\active \protected\def⇒{\ensuremath{\Rightarrow}}
